% =====================================================================
\section{Neural operators and Fourier layers}
\label{sec:neural-operators}
% =====================================================================

\subsection{Operator learning: problem setup}
\label{sec:setup}

A neural operator learns a map between fields rather than between vectors of a prescribed size. Let $\A$ and $\U$ be normed spaces of functions on domains $D$ and $D'$, and let $\G:\A\to\U$ denote the target map. Given input--output pairs $(f_j,g_j)$, with $g_j=\G(f_j)$, we seek an approximation $\Gtheta$ whose parameters can be used on different discretizations of those fields \citep{kovachki2023neural}. The domains need not be identical, although the architectures considered here predict on the input grid.

In this work, the input contains the matter-density field and, for the multi-field models, the LOS velocity. Redshift and simulation parameters specify which map is being emulated; they are conditioning variables, distinct from the learned weights $\theta$. The outputs are the neutral hydrogen fraction $\xHI$ and, where included, the brightness temperature $T_b$. The target is the map defined by the simulation and its preprocessing, rather than an exact solution operator of a specified PDE.

The function-space viewpoint separates an architecture from its numerical evaluation. It does not by itself guarantee continuity, resolution transfer or generalization. For example, $L^2$ is useful for field-error norms, whereas bounded measurable fields in $L^\infty$ are convenient for the gated interaction below. Point samples require a sufficiently regular representative; simulation cells can instead be understood as cell averages. Changing the grid must preserve the domain, physical units and meaning of every input channel.

\subsection{Neural operators as stacked kernel integrals}
\label{sec:neural-operator-layers}

A linear integral operator maps a field $f$ to a field $g$ through a kernel:
\begin{equation}
  \label{eq:kernel-integral}
  g(r)=\int_D\kappa(r,r';\theta)f(r')\dd r'.
\end{equation}
For vector-valued fields, $\kappa$ is a matrix that mixes channels. A Green's function gives one physical example of such a kernel, when the inverse of a linear differential operator admits an integral representation. A learned kernel need not be a Green's function or satisfy a differential equation.

An integral alone is linear when the kernel is fixed. A neural operator obtains a nonlinear map by combining integral interactions with pointwise transformations. A typical hidden layer is
\begin{equation}
  \label{eq:neural-operator-layer}
  v_{\ell+1}(r)=\sigma\!\left(W_\ell v_\ell(r)+b_\ell(r)
           +(\Kop_\ell v_\ell)(r)\right),
\end{equation}
with lifting $v_0=P(f)$ and projection $\Gtheta(f)=Q(v_L)$. Here $W_\ell$ mixes channels at one position, while $\Kop_\ell$ exchanges information between positions. This distinction matters: a fixed kernel does not make the complete network linear or its response independent of the input.

\subsection{The Fourier neural operator}
\label{sec:fno}

On a periodic domain, a translation-invariant kernel depends only on separation, $\kappa(r,r';\theta)=\kappa(r-r';\theta)$. Its integral is a convolution:
\begin{equation}
  \label{eq:kernel-integral-fourier}
  \Kop f=\Finv\!\left[\hat\kappa\,\F[f]\right].
\end{equation}
The Fourier neural operator (FNO) parameterizes the matrix-valued multiplier $\hat\kappa$ directly \citep{li2021fourier}. A lifting layer maps the input into $\dv$ channels, several layers of the form \cref{eq:neural-operator-layer} combine spectral and pointwise paths, and a projection produces the output (\cref{fig:fno-architecture}).

For a uniform grid, the integral is approximated by
\begin{equation}
  \label{eq:kernel-integral-discrete}
  g(r_i)=\sum_{j=1}^{N}\kappa(r_i-r_j;\theta)f(r_j)w_j,
\end{equation}
where $w_j$ is a cell-volume weight, not a one-dimensional spacing unless $d=1$. On the unit-volume domain used below, $w_j=1/N$. The FFT evaluates this discrete convolution without forming all $N^2$ pairs.

\begin{figure*}[!tp]
  \centering
  \figplaceholder[0.85\linewidth]{fno_architecture.pdf}{FNO block diagram: lift $P$ $\to$ Fourier layers ($\F$, $R_\phi$, $\Finv$ + local $W$) $\to$ projection $Q$.}
  \caption{Fourier neural operator architecture. The spectral path is linear; the pointwise activation makes the complete layer nonlinear.}
  \label{fig:fno-architecture}
\end{figure*}

\subsection{Discretization invariance and its limits}
\label{sec:discretisation}

Each spectral path retains a finite set of Fourier indices $\Lambda$:
\begin{equation}
  \label{eq:mode-truncation}
  \widehat{\Kop f}(k)=
  \begin{cases}
    R_\phi(k)\hat f(k), & k\in\Lambda,\\
    0, & k\notin\Lambda,
  \end{cases}
\end{equation}
with $R_\phi(k)\in\C^{\dv\times\dv}$. We use $\kmax$ for the mode-count setting and $\Lambda$ for the actual retained set. The implementation uses rectangular, possibly anisotropic bands rather than a spherical cutoff. On a box of side $L$, index $k$ corresponds to physical angular wavenumber $2\pi k/L$; a mode count is therefore not itself a wavenumber in $\mathrm{Mpc}^{-1}$.

The number of spectral parameters depends on $|\Lambda|$ and $\dv$, rather than on the grid size. Lifting, projection and pointwise maps also share their weights across cells. The weights can therefore be reused on finer grids of the same domain that resolve the retained modes. This is the architectural basis of the FNO's demonstrated zero-shot super-resolution \citep{li2021fourier}. Several qualifications matter here.

\paragraph{The spectral path is band-limited; the output need not be.} The pointwise path can retain high-frequency information already present in the hidden field, and nonlinear activations can generate additional frequencies. It is therefore incorrect to describe the complete FNO output as band-limited. A finer grid provides additional evaluation points, but reliable predictions on newly resolved physical scales still depend on what the model learned and on the information in its inputs.

\paragraph{Aliasing.} Nonlinear activations and products can create frequencies above the grid's Nyquist limit. Sampling folds these frequencies into the represented band, with an error that depends on the grid. Sharing parameters across resolutions is weaker than requiring the discrete evaluations to represent exactly the same continuous operator. This distinction is formalized by \citet{bartolucci2023representation}.

\paragraph{Resolution-dependent targets.} Changing the simulation resolution may change smoothing, sub-grid prescriptions and the resolved ionization fronts. Transferring between discretizations of one field and predicting a new simulation at finer resolution are different tests. An architecture that supports the former need not succeed at the latter.

\paragraph{Boundary conditions.} The FFT evaluates a periodic convolution. This matches coeval periodic boxes and the transverse directions of the lightcones. Along the non-periodic LOS, padding and cropping define a different bounded-domain operation. Padding reduces wrap-around effects only over the separations it actually separates; it does not impose a physical boundary condition by itself.

\subsection{Limitations of a stationary spectral interaction}
\label{sec:fno-limitations}

The FNO offers an efficient global interaction. For ionization fields, the questions are whether its retained modes and pointwise nonlinearities can represent sharp boundaries and whether the learned response adapts sufficiently to heterogeneous regions.

\paragraph{Stationarity of the linear kernel.} A spectral path applies the same linear kernel at every position. The full network can still have an input-dependent response: nonlinearities, conditioning and repeated layers change its Jacobian. What a standard spectral path lacks is an explicit source--receiver modulation of its kernel within that interaction. This is the distinction addressed by the rhizome.

\paragraph{Sharp features.} A discontinuity has a broad Fourier spectrum; for example, the coefficients of a one-dimensional step decay as $1/|k|$. Truncating a linear Fourier representation can smooth the edge and produce Gibbs ringing. Ionization maps are close to binary, so front placement is sensitive to these scales. Pointwise paths and activations can recover detail, but a small spectral band does not directly transport arbitrary high-frequency structure between cells.

\paragraph{Intermediate scales.} A standard FNO has no explicit spatial hierarchy or compact local stencil. It can represent interactions at intermediate distances through its kernel, so these scales are not absent; they simply have no dedicated branch as they do in a U-Net.

\paragraph{Cost of more modes.} A dense multiplier stores $\dv^2$ complex entries per retained mode. Its parameter count grows as $\mathcal{O}(|\Lambda|\dv^2)$, with $|\Lambda|$ growing cubically with an isotropic mode count in three dimensions. Factorizing the channel matrices can reduce this cost, but adds a restriction of its own.

These observations motivate a different interaction law; they do not establish an impossibility result for FNOs. The experiments must determine whether the proposed change improves front placement, cross-correlation and morphology under comparable inputs and training conditions.
